\section{Обсуждение ограничений и путей дальнейшего развития}

TODO2

На данный момент можно четко выделить два класса проблем, оказывающих влияние на точность, воспроизводимость и скорость развития системы.

Проблема качества и состава датасета. Датасет содержит неполные и неоднородные примеры, сильный дисбаланс по классам, чрезмерно сложные исходные модели. Имеющегося набора данных недостаточно, так как состав датасета по операциям неравномерен, часть примеров невалидна или непригодна для обучения, а истории построения часто используют не-атомарные операции, потому непригодные для обучения нейронной сети, которая должна делать отдельные простые шаги, собирая всё в итоге в дерево построения. Для некоторых операций из большого корпуса моделей удавалось использовать лишь небольшую долю наблюдений из-за отсутствия нужных шагов, пустых STL либо чрезмерной сложности исходных объектов. Решение состоит в жесткой валидации наборов, фильтрации нерепрезентативных примеров, выделении специализированных поддатасетов по операциям и формировании отдельного корпуса простых моделей для сегментации и параметризации. 

Проблема ограниченности STL как источника инженерной семантики. STL описывает поверхность как триангуляцию, но не содержит явной CAD-семантики граней, эскизов и ограничений. Решение — постепенный переход к комбинированному использованию STL для распознавания и STEP/B-Rep для уточнения инженерного смысла.


\begin{enumerate}
    \item Во время триангуляции, треугольники могут быть ориентированы по часовой или против часовой стрелки. На данный момент, ориентация определяется каким-то внутренним алгоритмом и, по сути, случайна. Можно использовать данный аспект как способ передать метаинформацию в модель - например, ориентировать треугольники по-разному в зависимости от типа порождающей поверхности, или ориентировать иначе треугольники на границе (стыке) двух поверхностей для облегчения детекции границ.
    \item Сейчас на инференсе STEP используется только для того, чтобы обеспечить известный алгоритм триангуляции, с фиксированным числом треугольников и рёбер. Если выбрать правильный метод конвертации (с равномерным заполнением точек), то можно любой входной файл конвертировать вначале в Point cloud, а потом в STL: STP -> Point cloud -> STL, STL -> Point cloud -> STL, …, обеспечив поддержку произвольного 3D формата. Таким образом, целесообразно добавить поддержку STL и/или Point cloud на инференсе путём преобразования всех данных (как на этапе обучения, так и на этапе инференса) - например, в Point cloud.
    \item Сейчас для обучения используется только последний шаг модели - берется геометрия, полученная на последнем шаге, и только она раскрашивается и используется при обучении. Так как триангуляция сильно меняется практически при любой операции, можно увеличить датасет используя, например, каждый второй шаг, или вообще каждый шаг построения модели как отдельную модель. Таким же образом можно будет использовать модели, у которых последний шаг слишком сложен, чтобы входить в обучающую выборку (слишком много треугольников) - использовать модель с первых N шагов, в которой треугольников становиться уже достаточно мало.
    \item Сейчас обучение происходит так, что предсказываются сразу все грани, порожденные выбранной операцией. Можно попробовать применить другой способ обучения, при котором обучение происходит только на последних последовательных однотипных операциях. Для примера возьмём cutExtrusion. Вместо того, чтобы искать сразу все вырезания, можно попробовать искать только все последние последовательные вырезания. Так, если будет 4 операции: bossExtrusion, cutExtrusion, bossExtrusion (который сделан поверх первого вырезания и сломавший его форму), а потом cutExtrusion - то мы не будем вводить штрафы для нейронной сети за то, что она не признала первое отверстие, которые было “сломано” выдавливанием, как мы бы делали в наивном подходе. Также, мы попутно решаем проблему “стола” - это если к столешнице были приделаны 4 ножки - то мы будем позволять bossExtrusion находить все 4 ножки сразу, независимо от того, в каком порядке их решил сделать исходный разработчик.
    \item Нейросеть фундаментально работает с рёбрами, и поэтому необходим алгоритм, преобразующий окрашенные рёбра в окрашенные треугольники. Сейчас в итоговый test\_0.stl попадают только те треугольники, которые содержат хотя бы 2 покрашенных ребра. Возможно, будет лучше выбирать только те, где покрашены все 3 ребра - хотя даже в таком случае (например, в случае тетраэдра) невозможно выбрать все грани кроме ровно одной. Возможно, стоит разбивать такие грани на более мелкие треугольники, и позволять нейросети таким образом выражать степень уверенности в каждой отдельной грани.
    \item ожно добавить в постобработке опцию для фильтрации треугольников не только по их площади (как происходит сейчас), а ещё, например, по длине, по углу наклона, по расстоянию от центра модели и прочим 3D характеристикам, которые получится найти - это позволит исключить ситуации, в которых нейронная сеть случайно выдала какой-то один длинный или дальний треугольник (например, находящийся сильно дальше от скучкованных всех других).
    \item Для всех операций можно обучить 3 версии нейросетей: одну для операции по изогнутым контурам, вторую для операции по контурам из отрезков, и третью для смешанных операций. Сейчас, так как число треугольников для этих трёх случаев сильно отличается, то нейронная сеть, обучающаяся на всех трёх случаях разом, стремится приоритизировать именно скруглённые контуры.
    \item В дальнейшем требуется дообучить остальные основные операции, такие как: rotated, linear\_pattern, axis3D, mirror\_pattern и пр.
    \item Можно применить скомбинированный принцип обучения: при наличии двух датасетов, один из которых большой, но “грязный”, а второй относительно маленький, но собранный и проверенный вручную, возможно обучить нейросеть вначале только на большом, чтобы она выучила основные паттерны и структуру, а потом уже дообучить на хорошем, чтобы она смогла разобраться в деталях и нюансах.
    \item Можно создать скрипт для автоматической оценки качества нейронной сети: ввести и зафиксировать конкретные метрики, такие как “процент разбирающихся в ноль моделей”, или “суммарное число шагов, на которое получилось упростить набор моделей”, или “общий объём удалённой части тела”, и запускать периодически его в связке с КОМПАС-3D для валидации и оценки новых вариантов и версий main.exe.
    \item Можно сделать отдельный программный модуль, который будет обучать и использовать нейронную сеть, которая при получении глобального набора треугольников "операции" будет выдавать границы между ними. Таким образом, мы сможем, получив вывод текущего сегментатора (например, что вся модель была сделана исключительно из bossExtrusion) вызвать этот “разделитель” и найти конкретно треугольники, относящиеся к каждому отдельному шагу.
    \item Наконец, можно увеличить количество треугольников в моделях, или размер самих нейронных сетей (с учетом увеличения вычислительных затрат), а также подобрать гиперпараметры обучения для улучшения качества моделей.
\end{enumerate}
